% BHU PYQ -- Manually transcribed & error-corrected
\documentclass[11pt,a4paper]{article}

\usepackage[a4paper,top=2.5cm,bottom=2.5cm,left=2.8cm,right=2.8cm,headheight=14pt]{geometry}
\usepackage{amsmath,amssymb}
\usepackage[T1]{fontenc}
\usepackage[utf8]{inputenc}
\usepackage{lmodern}
\usepackage{microtype}
\usepackage{fancyhdr}
\usepackage{enumitem}
\usepackage{hyperref}
\usepackage{lastpage}
\usepackage{parskip}

\hypersetup{colorlinks=true,urlcolor=black,linkcolor=black,
  pdftitle={CHB-301: Organic & Physical Chemistry-II -- BHU PYQ},pdfauthor={Banaras Hindu University}}

\pagestyle{fancy}\fancyhf{}
\renewcommand{\headrulewidth}{0.4pt}\renewcommand{\footrulewidth}{0.4pt}
\fancyhead[L]{\small   \textit{Banaras Hindu University}}
\fancyhead[R]{\small   \textit{CHB-301: Organic \& Physical Chemistry-II}}
\fancyfoot[C]{\small Page  \thepage\ of \pageref{LastPage}}


\newlist{parts}{enumerate}{2}
\setlist[parts,1]{label=(\alph*),leftmargin=*,itemsep=5pt,topsep=3pt}
\setlist[parts,2]{label=(\roman*),leftmargin=*,itemsep=3pt,topsep=2pt}

\newcommand{\pts}[1]{\hfill   \textit{[#1]}}


\begin{document}


\begin{center}
  {\large\bfseries BANARAS HINDU UNIVERSITY}\\[2pt]
  {
\normalsize B.Sc. (Hons.) Semester III Examination, 2022-23}\\[6pt]
  \rule{0.8\linewidth}{0.5pt}\\[6pt]
  {\large\bfseries Paper No. CHB-301: Organic \& Physical Chemistry-II}\\[4pt]
  {
\normalsize Time: 3 Hours\hfill Maximum Marks: 70}
\end{center}
\rule{\linewidth}{0.4pt}
\small



\noindent   \textit{Instructions: Answer five questions in total, including Question~1 from each section which is compulsory. Use separate answer books for Section~A and Section~B.}

\normalsize
\bigskip


\begin{center}
  {\large\bfseries SECTION A}\\[2pt]
  {
\normalsize (Organic Chemistry-II -- Marks: 35)}
\end{center}
\rule{0.4\linewidth}{0.2pt}
\smallskip




\noindent   \textbf{Question 1.} Answer all parts of the following: \pts{5 $ \times$ 2 = 10}

\begin{parts}
  \item Classify the following species as aromatic, antiaromatic, or nonaromatic: (i) Cyclopropenyl cation, (ii) Cyclopentadienyl anion, (iii) Cycloheptatriene, and (iv) Cyclooctatetraene.
  \item Among the following, which has the most acidic $\alpha$-hydrogen? Write the appropriate reason: (i) Acetone, (ii) Acetylacetone, (iii) Ethyl acetoacetate, and (iv) Acetaldehyde.
  \item Complete the following reactions, predicting the major product:
  
\begin{parts}
    \item Acetophenone $\xrightarrow{ \text{NaBH}_4} ?$
    \item Benzene $\xrightarrow{ \text{CH}_3 \text{COCl, AlCl}_3} ?$
  \end{parts}
  \item Write the main product in the following reaction, showing the mechanism:
    \[  \text{CF}_3 \text{-C}_6 \text{H}_4 \text{-Cl} \xrightarrow{ \text{NaNH}_2,  \text{NH}_3} ? \]
  \item Write the correct product of the following electrophilic aromatic substitution reaction with an appropriate reason:
    \[  \text{Toluene} \xrightarrow{ \text{HNO}_3,  \text{H}_2 \text{SO}_4} ? \]
\end{parts}

\medskip\rule{\linewidth}{0.2pt}




\noindent   \textbf{Question 2.}

\begin{parts}
  \item Draw the molecular orbital energy level diagram of benzene. How would you explain that all carbon-carbon bond lengths in benzene are identical? \pts{4}
  \item Compare the relative reactivity of nitrobenzene and phenol towards electrophilic aromatic substitution reactions. \pts{4}
  \item Write the mechanisms of the following reactions: \pts{4}
  
\begin{parts}
    \item Mannich reaction
    \item Acid-catalyzed aldol condensation
  \end{parts}
\end{parts}

\medskip\rule{\linewidth}{0.2pt}




\noindent   \textbf{Question 3.}

\begin{parts}
  \item What is positive and negative hyperconjugation? Explain with suitable examples. \pts{4}
  \item Write the mechanisms of the following rearrangement reactions: \pts{4}
  
\begin{parts}
    \item Hofmann rearrangement
    \item Haloform reaction
  \end{parts}
  \item Describe the DDT synthesis reaction mechanism and discuss its environmental uses and impacts. \pts{4}
\end{parts}

\medskip\rule{\linewidth}{0.2pt}




\noindent   \textbf{Question 4.}

\begin{parts}
  \item Discuss the Reimer-Tiemann reaction with a detailed stepwise mechanism. \pts{4}
  \item Compare and explain the relative reactivity of acid chloride, acid anhydride, amide, and ester with a nucleophile. \pts{4}
  \item Explain the mechanism for the synthesis of diazonium salts. How would you obtain the following from benzenediazonium chloride? \pts{4}
  
\begin{parts}
    \item Thiophenol
    \item Phenyl isocyanide
    \item Biphenyl
    \item Phenylhydrazine
  \end{parts}
\end{parts}


\newpage

\begin{center}
  {\large\bfseries SECTION B}\\[2pt]
  {
\normalsize (Physical Chemistry-II -- Marks: 35)}
\end{center}
\rule{0.4\linewidth}{0.2pt}
\smallskip




\noindent   \textbf{Question 5.} Answer all parts of the following: \pts{5 $ \times$ 2 = 10}

\begin{parts}
  \item What is liquid junction potential? How can it be minimized?
  \item Compare the transference numbers of the chloride ion ($ \text{Cl}^-$) in $ \text{HCl}$ and $ \text{NaCl}$ solutions.
  \item The resistance of a $0.5\, \text{M}$ solution of an electrolyte in a cell was found to be $45\, \text{ohm}$. Calculate the molar conductance of the solution if the electrodes of the cell are $2.2\, \text{cm}$ apart and have an area of $3.8\, \text{cm}^2$.
  \item Calculate the Miller indices of the crystal planes for which the intercepts along the crystallographic axes are $(a/2, 2b/3, \infty c)$.
  \item Write the Gibbs phase rule and define the terms involved taking some suitable examples.
\end{parts}

\medskip\rule{\linewidth}{0.2pt}




\noindent   \textbf{Question 6.}

\begin{parts}
  \item How is the molar conductivity at infinite dilution ($\Lambda_m^\infty$) determined for acetic acid ($ \text{CH}_3 \text{COOH}$) and hydrochloric acid ($ \text{HCl}$)? Explain using Kohlrausch's law. \pts{6}
  \item What would be the pH of a solution in a cell if it generates $0.833\, \text{V}$ e.m.f. at $25^\circ \text{C}$, when it is connected to a saturated calomel electrode (SCE)? \pts{4}
  \item What do you mean by a reversible electrode? Discuss any one reversible electrode mentioning its Nernst equation for electrode potential. \pts{4}
\end{parts}

\medskip\rule{\linewidth}{0.2pt}




\noindent   \textbf{Question 7.}

\begin{parts}
  \item Discuss the determination of transport numbers using the moving boundary method. \pts{6}
  \item What is the degree of freedom for a system consisting of a saturated aqueous solution of $ \text{NaCl}$ in contact with solid $ \text{NaCl}$ and its vapours? \pts{4}
  \item Draw and discuss the phase diagram of the lead-silver ($ \text{Pb-Ag}$) system. What is the eutectic point? \pts{4}
\end{parts}

\medskip\rule{\linewidth}{0.2pt}




\noindent   \textbf{Question 8.}

\begin{parts}
  \item What is the Nernst distribution law? Describe the procedure to determine the equilibrium constant of the reaction $ \text{KI} +  \text{I}_2 \rightleftharpoons  \text{KI}_3$ in aqueous medium using the Nernst distribution law. \pts{6}
  \item Define space lattice and unit cell of a crystal. \pts{2}
  \item Write a short note on different types of point defects (Schottky and Frenkel defects) in solids. \pts{3}
  \item Derive Bragg's equation for X-ray diffraction. What is the significance of $n$ in this equation? \pts{3}
\end{parts}

\vfill
\rule{\linewidth}{0.4pt}

\begin{center}\small
  \href{https://bhu-exam.vercel.app/}{Click here to visit our website}\\[4pt]
  \href{https://me-aryan.vercel.app/}{Click here to visit our contributor}\\[4pt]
  \href{https://quashan-me.vercel.app/}{Click here to visit our contributor}
\end{center}

\end{document}
